\section{נתונים, תכונות ופרוטוקול הניסוי}

\subsection{פאנל הנכסים}

כדי לבדוק שהמסקנות אינן תלויות במניה אחת בלבד, נבחר פאנל של תשעה נכסים בעלי מאפיינים שונים. המטרה אינה לבצע בחירת מניות, אלא להפעיל אותה מתודולוגיה על סוגי סיכון שונים ולבדוק האם התנהגות המודל נשארת דומה.

הנתונים נמשכו באמצעות \tech{yfinance} \cite{yfinance} עבור חלון בקשה של \datecode{2014-01-01 - 2024-12-31}. הנכסים הם:

\begin{itemize}
    \item \asset{SPY} - קרן סל המייצגת את שוק המניות האמריקאי הרחב. היא משמשת גם כנכס הייחוס בניתוח הרב־נכסי.
    \item \asset{QQQ} - קרן סל עם משקל גבוה לחברות טכנולוגיה וצמיחה גדולות.
    \item \asset{IWM} - קרן סל המייצגת מניות אמריקאיות קטנות יחסית.
    \item \asset{TLT} - אג"ח ממשלתיות אמריקאיות ארוכות טווח.
    \item \asset{GLD} - זהב.
    \item \asset{HYG} - אג"ח חברות בדירוג גבוה־סיכון יחסית, כלומר \tech{high yield}.
    \item \asset{BTC-USD} - ביטקוין, כנכס אלטרנטיבי בעל תנודתיות גבוהה.
    \item \asset{JPM} - מניה גדולה מסקטור הפיננסים.
    \item \asset{NVDA} - מניית צמיחה בעלת רגישות גבוהה יחסית לשוק.
\end{itemize}

חשוב להבחין בין חלון הבקשה לבין הנתונים שנכנסו בפועל למודל. זמינות הנתונים שונה מעט בין הנכסים, ובפרט סדרת \asset{BTC-USD} מתחילה מאוחר יותר. לכן כל חלוקה ל־\tech{Train}, \tech{Validation} ו־\tech{Test} נעשתה מתוך התצפיות הזמינות בפועל עבור אותו נכס.

ההורדה בוצעה עם \tech{auto\_adjust=False}. כאשר העמודה \tech{Adj Close} הייתה זמינה, היא שימשה לחישובי תשואה ויעד; אחרת נעשה שימוש ב־\tech{Close}. לא בוצעה השלמה מלאכותית של ימים חסרים. שורות עם ערכי \tech{OHLCV} חסרים או לא־סופיים הוסרו, ולאחר יצירת התכונות הוסרו גם שורות שאין להן מספיק היסטוריה לחישוב חלון מתגלגל או שאין עבורן יעד של היום הבא. בניתוח בין נכסים השתמשנו רק בתאריכים שבהם קיימת תצפית בשתי הסדרות המושוות.

\subsection{וקטור התכונות}

בכל יום המודל מקבל תיאור קצר של מצב השוק באותו יום. במקום להזין עשרות אינדיקטורים טכניים, נבחרו ארבע תכונות פשוטות שקל לפרש. הן מנסות לתאר ארבעה היבטים שונים: כיוון התנועה, עוצמת התנודתיות, טווח המסחר בתוך היום ופעילות המסחר.

וקטור התצפית בזמן \(t\) הוא:

\begin{equation}
\mathbf{x}_t = \left[r_t,\; \sigma_t^{(20)},\; R_t,\; \Delta\log V_t\right].
\end{equation}

לכל רכיב יש משמעות ישירה:

\begin{itemize}
    \item \textbf{תשואה לוגריתמית, \(r_t\).} מודדת את שינוי המחיר בין יום \(t-1\) ליום \(t\):
    \begin{equation}
    r_t = \log\left(\frac{P_t}{P_{t-1}}\right).
    \end{equation}
    ערך חיובי מצביע על עלייה וערך שלילי על ירידה.

    \item \textbf{תנודתיות מתגלגלת, \(\sigma_t^{(20)}\).} סטיית התקן של התשואות ב־20 הימים האחרונים. תכונה זו מתארת אם התקופה האחרונה הייתה רגועה או תנודתית.

    \item \textbf{טווח יומי, \(R_t\).} מבוסס על המחיר הגבוה והנמוך של אותו יום. הוא נותן מידע על עוצמת התנועה בתוך יום המסחר, גם אם מחיר הסגירה השתנה מעט.

    \item \textbf{שינוי לוגריתמי בנפח, \(\Delta\log V_t\).} מודד האם נפח המסחר עלה או ירד ביחס ליום הקודם.
\end{itemize}

כל הנתונים שמשמשים לחישוב התכונות של יום \(t\) מגיעים מהעבר ומהיום הנוכחי בלבד. אין שימוש במחיר או בנפח של יום \(t+1\) בעת בניית וקטור הקלט של יום \(t\).

\subsection{יעד החיזוי והחלוקה הכרונולוגית}

משימת החיזוי היא בינארית: בכל יום אנו שואלים האם התשואה של יום המסחר הבא תהיה חיובית או שלילית. כלומר, הקלט מגיע מיום \(t\), ואילו התווית מתייחסת ליום \(t+1\).

היעד מוגדר כך:

\begin{equation}
y_t = \mathbb{1}\left[r_{t+1}\ge 0\right].
\end{equation}

אם התשואה ביום הבא אינה שלילית, מתקבל \(y_t=1\); אחרת מתקבל \(y_t=0\).

בגלל שמדובר בסדרת זמן, אסור לערבב את הדוגמאות באקראי. חלוקה אקראית הייתה עלולה לאפשר למודל ללמוד מתקופה מאוחרת ולבדוק את עצמו על תקופה מוקדמת יותר, מצב שאינו אפשרי בעולם אמיתי. לכן החלוקה היא כרונולוגית, בקירוב:

\begin{itemize}
    \item \pct{70} ראשונים - \tech{Train}, לצורך התאמת המודל.
    \item \pct{15} הבאים - \tech{Validation}, לצורך בחירת תצורת ה־HMM.
    \item \pct{15} אחרונים - \tech{Test}, לצורך הערכה סופית בלבד.
\end{itemize}

לא בוצע \tech{shuffle}. בנוסף, כדי למנוע זליגה בדיוק בגבולות החלוקה, השורה האחרונה של כל חלק מוסרת אם היעד שלה כבר שייך לחלק הבא. כך אף תווית עתידית אינה חוצה את הגבול בין קבוצות הנתונים.

יש להבחין בין שלב בחירת המודל לבין ההתאמה הסופית. בזמן השוואת מועמדי HMM, ה־\tech{scaler} מותאם על \tech{Train} בלבד, כל מועמד מאומן על \tech{Train}, והביצועים שלו נמדדים על \tech{Validation}. רק לאחר שנבחרו מספר המצבים \(K\) וה־\tech{seed}, המודל הנבחר וה־\tech{scaler} מותאמים מחדש על כל המידע שהיה זמין לפני ה־\tech{Test}, כלומר \tech{Train+Validation}. ה־\tech{Test} עצמו נשאר סגור עד שלב ההערכה הסופי.

במודלי \tech{Supervised Learning} נקבעו ה־\tech{hyperparameters} מראש ולא בוצעה בחירה לפי \tech{Validation}. לכן שלושת המודלים האלה אומנו על \tech{Train} בלבד ונבדקו על אותו \tech{Test} כרונולוגי. המשמעות היא שההשוואה בינם לבין ה־HMM אינה משתמשת בדיוק באותה כמות היסטוריה לאימון; נקודה זו מצוינת בהמשך כחלק ממגבלות הניסוי.

\subsection{שחזור ומעקב אחר מקור התוצאות}

כל ריצה נשמרת בתיקייה חדשה יחד עם קובץ \tech{manifest}. הקובץ מתעד את התאריך, גרסת הקוד, גרסאות הספריות, רשימת הנכסים, התכונות, האתחולים והגדרות המודל. כך ניתן לקשר כל מספר בדוח לריצה שממנה הוא הגיע.

ריצת ה־HMM המרכזית היא \code{extended_20260811_121957}, וריצת מודלי ה־\tech{Supervised Learning} היא \code{supervised_20260811_133344}. תוצאות ישנות נשמרו לצורכי ביקורת ושחזור, אך אינן משמשות לטענות המספריות בדוח זה.
